\section{Introduction}
\label{sec:intro}

Modern decision systems learn from logs generated by earlier versions of themselves. Recommendation, advertising, marketplace ranking, experimentation, and clinical decision support all use a logging policy to select an exposure or action, observe only the corresponding outcome, and then update a model or policy. The scientific questions asked of those logs are broader than average policy value. Practitioners also need to know whether a score remains calibrated under an intervention, whether a target policy violates a safety or resource constraint, whether treatment effects differ across hidden subpopulations, and whether exposure produces unequal direct or spillover effects. Each question can be written as a \emph{counterfactual moment}: an expectation of an intervention-specific orthogonal score multiplied by an auditor or constraint function.

Reliable estimation of these moments is difficult when three pathologies interact. First, overlap may be weak. A target policy can place mass on actions that the logging policy chose rarely, producing large importance ratios even when outcomes are bounded. Strong overlap assumptions hide this difficulty and can fail precisely for the policies that are most useful to evaluate. Second, the outcome itself may be heavy-tailed: watch time, transaction value, medical cost, and user-level aggregate engagement are not naturally bounded. Third, logged units are rarely independent. Candidates share a slate, repeated observations share users, neighboring nodes interfere, and policy updates induce temporal dependence. Treating such data as independent can yield narrow intervals with severe undercoverage.

Existing solutions address important pieces of this problem. Doubly robust estimators reduce sensitivity to nuisance estimation \citep{dudik2011doubly,jiang2016doubly,farajtabar2018more}; weak-overlap analyses characterize when truncation is necessary \citep{wang2017optimal,mehrabi2024weak,khan2024beyond}; robust mean estimators tolerate finite moments \citep{lugosi2019mean}; and causal inference under interference formalizes exposure mappings and dependency structures \citep{aronow2017estimating,savje2021average}. What is missing is a common interface that simultaneously supports a \emph{family} of counterfactual questions, learned nuisance functions, heavy-tailed scores, and observable graph--time dependence---and then turns the resulting uncertainty into decisions rather than a post-hoc error bar.

We introduce \textbf{\COMET} (\textbf{CO}unterfactual \textbf{M}oment \textbf{E}stimation and \textbf{T}esting). \COMET has a modular structure. A generalized exposure mapping converts a possibly joint action into a local exposure state. A cross-fitted orthogonal score converts the partially observed outcome into a pseudo-outcome for a target intervention. A graph--time separation step prevents nuisance leakage and records which scores can be treated as independent or weakly coupled. Finally, a heavy-tail robust aggregation module produces simultaneous confidence radii over a policy family and an auditor family. The same engine can estimate policy value, certify counterfactual calibration, compare subgroup effects, test spillovers, or enforce a cost constraint.

The statistical engine supports a broader workflow, illustrated in \cref{fig:overview}. On a discovery split, a flexible auditor searches for a subgroup with a large counterfactual residual. The selected auditor is frozen and evaluated on a dependence-separated certification split, avoiding the double use of data that otherwise turns subgroup search into an invalid significance test. The resulting upper and lower confidence bounds can then be used in a pessimistic policy-selection problem: utility is optimized through a lower confidence bound while every risk or fairness constraint is enforced through an upper confidence bound. We refer to this workflow as \emph{discover--certify--repair}.

\begin{figure}[t]
  \centering
  \resizebox{\linewidth}{!}{\begin{tikzpicture}[
  node distance=5mm and 5mm,
  box/.style={draw,rounded corners=1.5pt,align=center,minimum height=8mm,inner xsep=5pt,inner ysep=3pt,font=\small},
  arrow/.style={-{Latex[length=2mm]},semithick},
  group/.style={draw,dashed,rounded corners=2pt,inner sep=4pt}
]
\node[box] (log) {Adaptive logs\\$X,A,S,Y,e$};
\node[box,right=of log] (score) {Orthogonal score\\$\psi^q(\widehat\eta)$};
\node[box,right=of score] (split) {Graph--time\\separation};
\node[box,right=of split] (robust) {Robust aggregation\\and radii};
\draw[arrow] (log) -- (score);
\draw[arrow] (score) -- (split);
\draw[arrow] (split) -- (robust);

\node[box,below=7mm of score] (discover) {Discover\\learn $h$ \\ on split 1};
\node[box,below=7mm of split] (certify) {Certify\\test frozen \\ $h$ on split 2};
\node[box,below=7mm of robust] (repair) {Repair\\pessimistic\\ policy choice};
\draw[arrow] (score) -- (discover);
\draw[arrow] (discover) -- (certify);
\draw[arrow] (robust) -- (certify);
\draw[arrow] (certify) -- (repair);

\node[group,fit=(score)(split)(robust),label={[font=\scriptsize]above:counterfactual moment engine}] {};
\end{tikzpicture}
}
  \caption{\COMET separates the counterfactual score, the dependence-aware robust estimator, and downstream use. Rich auditors may be learned on a discovery split, but certification and policy repair use a separate graph--time split.}
  \label{fig:overview}
\end{figure}

Our formulation is deliberately not tied to a single fairness metric or graph task. Counterfactual online performative fairness on evolving graphs \citep{li2026copf} is recovered by choosing a shown/not-shown exposure mapping, group--score auditors, and exposure-effect constraints. In the present work, this is one application rather than the organizing premise. The generalization changes both the assumptions and the scientific question: strict overlap is replaced by finite-moment weak overlap; fixed hand-designed slices are augmented by learned auditors; a stationary population mean is replaced by a window target with explicit drift; and a generic stochastic optimizer is replaced by certificate-guided policy repair.

\paragraph{Contributions.}
The paper develops the following components.
\begin{enumerate}[leftmargin=1.4em,itemsep=1pt,topsep=2pt]
    \item \textbf{A counterfactual-moment interface.} We give one orthogonal-score formulation that covers off-policy value, subgroup value, counterfactual calibration, treatment-effect heterogeneity, safety and resource constraints, and local spillovers.
    \item \textbf{Dependence-aware robust certification.} We combine graph--time separated cross-fitting with robust block aggregation. The resulting target bound decomposes into a finite-$p$ heavy-tail term with rate $N_{\eff}^{-(p-1)/p}$, a nuisance-product remainder, a temporal coupling term, and a within-window drift term. The dependence penalty is based on a computable conflict graph rather than an unobserved statement that samples are ``approximately i.i.d.''
    \item \textbf{Discover--then--certify auditing.} A rich model can search for a violating subgroup on one split; a frozen auditor is then certified on another. Ratio metrics include uncertainty in both the numerator and subgroup mass and abstain when support is insufficient.
    \item \textbf{Certificate-guided policy repair.} We formulate pessimistic constrained policy selection and show that simultaneous confidence bounds imply high-probability feasibility for the selected policy.
    \item \textbf{A falsifiable empirical program.} Experiments are designed around coverage, interval width, false certification, detection power, effective-sample-size scaling, subgroup recovery, and utility--constraint tradeoffs on synthetic, semi-synthetic, and real logged data. No conclusion relies only on point-estimation error.
\end{enumerate}

\draftnote{The strongest ICLR version should include: (i) a complete dependence-aware upper bound, (ii) either a matching lower bound or an adaptive unknown-tail result, and (iii) experiments with empirical coverage. The current bundle states the intended theorem precisely enough to guide those proofs but does not treat unfinished proof steps or placeholder tables as completed results.}
